\documentclass[11pt,a4paper]{article}
\usepackage[left=24mm,right=24mm,top=22mm,bottom=20mm,includefoot]{geometry}
\usepackage{amsmath,booktabs}
\usepackage{iftex}
\ifPDFTeX
  \usepackage{CJKutf8}
  % CJKutf8 supplies its own Unicode map for each Chinese subfont.
  \makeatletter
  \AddToHook{cmd/CJK@addcmap/before}{\pdfnobuiltintounicode\font@name}
  \makeatother
\else
  \usepackage[no-math]{fontspec}
  \IfFontExistsTF{Noto Serif CJK SC}{
    \setmainfont{Noto Serif CJK SC}
    \setsansfont{Noto Sans CJK SC}
  }{
    \setmainfont{FandolSong-Regular.otf}[BoldFont=FandolSong-Bold.otf]
    \setsansfont{FandolHei-Regular.otf}[BoldFont=FandolHei-Bold.otf]
  }
  \ifXeTeX
    \XeTeXlinebreaklocale "zh"
    \XeTeXlinebreakskip=0pt plus 0.5pt
  \fi
\fi
\pagestyle{plain}
\raggedbottom
\linespread{1.13}
\setlength{\parindent}{2em}
\setlength{\parskip}{3pt plus 1pt}
\setlength{\jot}{3pt}
\DeclareMathSizes{10.95}{10.95}{9}{7}
% Use a single stroke for Chinese headings when pdfLaTeX has no bold CJK font.
\newcommand{\headingweight}{\ifPDFTeX\mdseries\else\bfseries\fi}
\makeatletter
\renewcommand{\section}{\@startsection{section}{1}{0pt}
  {2.0ex plus .3ex minus .2ex}{.9ex}{\normalfont\large\headingweight}}
\makeatother

\begin{document}
\ifPDFTeX\begin{CJK*}{UTF8}{gbsn}\fi
\begin{center}
  {\LARGE\headingweight 两个暴露对有序结局的 Shapley 贡献}
\end{center}

设结局 $Y$ 分为一级、二级和三级，依次记为 $1,2,3$，其中三级最严重。
按四个潜在结果 $(Y_{00},Y_{01},Y_{10},Y_{11})$ 的顺序记录个体的类型
$t=(t_{00},t_{01},t_{10},t_{11})$。例如，$t=1123$ 表示
$Y_{00}=1,\ Y_{01}=1,\ Y_{10}=2,\ Y_{11}=3$。

记 $\mathcal E_k=(X_1=1,X_2=1,Y=k)$。由一致性，在 $\mathcal E_3$ 下有 $Y_{11}=3$。
每升一个等级计 $1$ 分。对给定类型 $t$，用 $R(j\mid t,\mathcal E_3)$ 表示暴露 $X_j$
对结局从 $Y_{00}$ 升至三级的 Shapley 贡献，单位为分。
给定类型 $t$，四个潜在结果取该类型的值，即 $Y_{ab}=t_{ab}$。两个暴露的 Shapley 公式为
\[
\begin{aligned}
 R(1\mid t,\mathcal E_3)
   &=\tfrac12\big[(Y_{10}-Y_{00})+(Y_{11}-Y_{01})\big],\\
 R(2\mid t,\mathcal E_3)
   &=\tfrac12\big[(Y_{01}-Y_{00})+(Y_{11}-Y_{10})\big].
\end{aligned}
\]
每一行的第一项是该暴露先加入时的等级增量，第二项是它后加入时的等级增量。
两种顺序各占一半。两个暴露的贡献之和为 $Y_{11}-Y_{00}=3-Y_{00}$。

\begin{center}
\renewcommand{\arraystretch}{1.06}
\setlength{\tabcolsep}{9pt}
\begin{tabular}{@{}l c c c@{}}
\toprule
观测等级 $k$ & 类型 $t$ & $R(1\mid t,\mathcal E_k)$ & $R(2\mid t,\mathcal E_k)$\\
\midrule
三级 & $1113$ & $1$ & $1$ \\
 & $1123$ & $3/2$ & $1/2$ \\
 & $1133$ & $2$ & $0$ \\
 & $1213$ & $1/2$ & $3/2$ \\
 & $1223$ & $1$ & $1$ \\
 & $1233$ & $3/2$ & $1/2$ \\
 & $1313$ & $0$ & $2$ \\
 & $1323$ & $1/2$ & $3/2$ \\
 & $1333$ & $1$ & $1$ \\
 & $2223$ & $1/2$ & $1/2$ \\
 & $2233$ & $1$ & $0$ \\
 & $2323$ & $0$ & $1$ \\
 & $2333$ & $1/2$ & $1/2$ \\
 & $3333$ & $0$ & $0$ \\
\midrule
二级 & $1112$ & $1/2$ & $1/2$ \\
 & $1122$ & $1$ & $0$ \\
 & $1212$ & $0$ & $1$ \\
 & $1222$ & $1/2$ & $1/2$ \\
 & $2222$ & $0$ & $0$ \\
\midrule
一级 & $1111$ & $0$ & $0$ \\
\bottomrule
\end{tabular}
\end{center}
\noindent 表中列出单调性条件 $Y_{00}\le Y_{01}\le Y_{11}$ 和 $Y_{00}\le Y_{10}\le Y_{11}$ 下的全部类型，
按双暴露时的观测等级 $k=Y_{11}$ 分组。责任值的单位为分。

\ifPDFTeX\end{CJK*}\fi
\end{document}
